\makeabstract
{
Avoided-Crossing Coupling and Control of DPPH Spins On-Chip at Cryogenic Temperatures
}
{
Stellan Small (1), Mridul Pushp (2), Jonathan Hess (3), Anthony Sigillito, PhD (2)
}
{
\textsuperscript{1}Department of Physics, Amherst College, Amherst, MA, USA\\
\textsuperscript{2}Department of Electrical and Systems Engineering, University of Pennsylvania, Philadelphia, PA, USA\\
\textsuperscript{3}Department of Physics, University of Pennsylvania, Philadelphia, PA, USA
}
{
Silicon spin quantum dots are a promising quantum computing paradigm due to their scalability and compatibility with widespread industrial fabrication techniques. The standard method of controlling these quantum dots is through pulsed microwave signals, but integrating these signals into the quantum dot environment poses significant challenges. Coplanar waveguides (CPWs) are a chip architecture that allows electrical components to be written onto the same chip as a silicon spin quantum dot, providing a path to interfacing a quantum dot with a larger apparatus. Additionally, prior research has shown the quantum effect of avoided crossing to be a viable means of coupling a spin system to a power feed line. We aim to create a CPW component with the potential to couple to silicon spin quantum dots via avoided crossing and integrate with other components in a quantum computational apparatus, such as a superconducting Mach-Zehnder amplitude modulator (AM). 
}
{
We designed a component where an electrical feed line is capacitively coupled to hangar resonators of various sizes which connect to the CPW ground plane at one end and etched this design into Niobium deposited onto a Silicon wafer. We tested avoided crossing capabilities by applying DPPH, a standard material in electron spin experiments and therefore a good substitute for quantum dots. We applied DPPH to the resonators in two ways: suspended in isopropanol or in vacuum grease
}
{
From measuring the frequency spectrum of the circuit as a function of magnetic field, we found that all four resonators which had vacuum grease applied to them experienced a noticeable effect when the cavity and DPPH spin systems were on-resonance, but only two of them displayed the expected mode-splitting behavior caused by avoided crossing. Further testing showed that one of the fully coupling cavities is within the 100 MHz tunability range of a mode of an AM connected in series.
}
{
. These results suggest the viability of CPW hangar resonators for quantum control applications and allow our group to proceed with testing the amplitude modulation component, which advances our long-term goal of producing low-dissipation quantum circuitry components.
}

\newpage
